\documentclass[12pt]{article}

\usepackage{amsmath, amssymb}
\usepackage{geometry}
\geometry{margin=1in}

\title{\textbf{CSE400 -- Lecture 7}\\
Expectation, CDFs, PDFs and Problem Solving}
\date{January 27, 2025}

\begin{document}

\maketitle

\section*{1. Lecture Structure (Page 2)}

According to the outline:

\begin{itemize}
\item Cumulative Distribution Function (CDF)
\begin{itemize}
\item Definition
\item Properties
\item Examples
\end{itemize}

\item Probability Density Function (PDF)
\begin{itemize}
\item Definition
\item CDF--PDF Relationship
\end{itemize}

\item Expectation of Random Variables (introduced later)
\end{itemize}

The lecture follows this order.

\section*{2. Intuition: CDF and PDF (Water Tank Analogy)}

\subsection*{Given Model}

A cylindrical water tank:

\begin{itemize}
\item Radius = $R$
\item Height = $H$
\item Water height = $h$
\end{itemize}

\subsection*{Interpretation}

The slide states:

\begin{quote}
Height $h \rightarrow$ value of random variable $X$\\
Volume up to $h \rightarrow$ cumulative probability
\end{quote}

So:

\begin{itemize}
\item Each height corresponds to a value of $X$
\item Accumulated water corresponds to $\Pr(X \le h)$
\end{itemize}

Thus:

\[
\text{Volume up to } h \;\longleftrightarrow\; F_X(h)
\]

\subsection*{Volume Calculation}

Given:

\[
V(h) = \int_{0}^{h} \pi R^2 \, dh
\]

Step-by-step:

\begin{itemize}
\item Cross-sectional area = $\pi R^2$
\item Volume = Area $\times$ Height
\item Integrate area from $0$ to $h$
\end{itemize}

Result:

\[
V(h) = (\pi R^2)h
\]

\subsection*{Probability Interpretation}

Maximum volume:

\[
V(H) = \pi R^2 H
\]

This corresponds to total probability = $1$.

Thus:

\begin{itemize}
\item $\pi R^2$ behaves like a constant PDF
\item $V(h)$ behaves like a CDF
\end{itemize}

CDF = accumulated PDF.

\section*{3. Cumulative Distribution Function (CDF)}

\subsection*{3.1 Definition}

The slide defines:

\[
F_X(x) = \Pr(X \le x), \quad -\infty < x < \infty
\]

Step-by-step meaning:

\begin{itemize}
\item Fix a real number $x$
\item Consider the event: ``$X \le x$''
\item Compute its probability
\item That probability is $F_X(x)$
\end{itemize}

So:

The CDF gives the total probability accumulated up to $x$.

\subsection*{Importance}

The slide states:

Most information about the random experiment is determined by the behavior of $F_X(x)$.

Meaning:

Knowing $F_X(x)$ allows us to compute all probabilities.

CDF fully characterizes the random variable.

\subsection*{3.2 Properties of CDF}

\subsubsection*{Property 1: Boundedness}

\[
0 \le F_X(x) \le 1
\]

Reasoning:

$F_X(x)$ is a probability.

\subsubsection*{Property 2: Limits}

\[
F_X(-\infty)=0, \quad F_X(\infty)=1
\]

Reasoning:

\begin{itemize}
\item At $-\infty$, probability is $0$
\item At $+\infty$, probability is $1$
\end{itemize}

\subsubsection*{Property 3: Monotonicity}

For $x_1 < x_2$:

\[
F_X(x_1) \le F_X(x_2)
\]

Reasoning:

\[
\{X \le x_1\} \subset \{X \le x_2\}
\]

\subsubsection*{Property 4: Interval Probability}

For $x_1 < x_2$:

\[
\Pr(x_1 < X \le x_2)=F_X(x_2)-F_X(x_1)
\]

Derivation:

\[
\Pr(X \le x_2)=\Pr(X \le x_1)+\Pr(x_1<X\le x_2)
\]

Rearranging:

\[
\Pr(x_1<X\le x_2)=F_X(x_2)-F_X(x_1)
\]

\section*{4. CDF Example \#1: Validity Check}

\subsection*{(a)}

\[
F_X(x)=\frac{1}{2}+\frac{1}{\pi}\tan^{-1}(x)
\]

Range of $\tan^{-1}(x)$ is $(-\pi/2,\pi/2)$.

So range is $(0,1)$.

Valid CDF.

\subsection*{(b)}

\[
F_X(x)=(1-e^{-x})u(x)
\]

Valid CDF.

\subsection*{(c)}

\[
F_X(x)=e^{-x^2}
\]

Invalid.

\subsection*{(d)}

\[
F_X(x)=x^2u(x)
\]

Invalid.

\section*{5. CDF Example \#2}

Given:

\[
F_X(x)=(1-e^{-x})u(x)
\]

\subsection*{(a)}

\[
\Pr(X>5)=1-F_X(5)=e^{-5}
\]

\subsection*{(b)}

\[
\Pr(X<5)=F_X(5)=1-e^{-5}
\]

\subsection*{(c)}

\[
\Pr(3<X<7)=F_X(7)-F_X(3)=e^{-3}-e^{-7}
\]

\subsection*{(d)}

Let $A=\{X>5\}$, $B=\{X<7\}$.

\[
\Pr(A|B)=\frac{F(7)-F(5)}{F(7)}
\]

\section*{6. Probability Density Function (PDF)}

\subsection*{6.1 Definition}

\[
f_X(x)=\lim_{\varepsilon\to 0}
\frac{\Pr(x<X<x+\varepsilon)}{\varepsilon}
\]

\subsection*{6.2 Using CDF}

\[
\Pr(x<X<x+\varepsilon)=F_X(x+\varepsilon)-F_X(x)
\]

\[
f_X(x)=\lim_{\varepsilon\to 0}
\frac{F_X(x+\varepsilon)-F_X(x)}{\varepsilon}
\]

\subsection*{6.3 PDF--CDF Relationship}

\[
f_X(x)=\frac{dF_X(x)}{dx}
\]

\subsection*{6.4 Inverse Relationship}

\[
F_X(x)=\int_{-\infty}^{x} f_X(t)\,dt
\]

\section*{7. Expectation of Random Variables}

Only heading is visible.

No further content available.

\section*{Final Exam-Oriented Logical Map}

\subsection*{CDF}

\[
F_X(x)=\Pr(X\le x)
\]

Properties:

\begin{itemize}
\item $0\le F\le 1$
\item $F(-\infty)=0,\;F(\infty)=1$
\item Non-decreasing
\item $\Pr(a<X\le b)=F(b)-F(a)$
\end{itemize}

\subsection*{PDF}

\[
f(x)=\lim_{\varepsilon\to 0}
\frac{\Pr(x<X<x+\varepsilon)}{\varepsilon}
\]

\[
f(x)=\frac{dF(x)}{dx}
\]

\[
F(x)=\int_{-\infty}^{x} f(t)\,dt
\]

\end{document}

